% ============ §3 Chain conditions (原书 pp.14-19) ============
\section{Chain conditions}\label{sec:3}\index{Chain conditions, section 3}

The following two conditions on a partially ordered set $\Gamma$ are
equivalent:
\begin{itemize}
  \item[(*)] any non-empty subset of $\Gamma$ has a maximal element;
  \item[(**)] any ascending chain $\gamma_1 < \gamma_2 < \cdots$ of elements of $\Gamma$ must stop after a
        finite number of steps.
\end{itemize}
The implication $(*) \Rightarrow (**)$ is obvious. We prove $(**) \Rightarrow (*)$. Let $\Gamma'$ be a non-
empty subset of $\Gamma$. If $\Gamma'$ does not have a maximal element, then by the axiom
of choice, for each $\gamma \in \Gamma'$ we can choose a bigger element of $\Gamma'$, say $\varphi(\gamma)$. Now if
we choose any $\gamma_1 \in \Gamma'$ and set $\gamma_2 = \varphi(\gamma_1)$, $\gamma_3 = \varphi(\gamma_2), \ldots$ then we get an infinite
ascending chain $\gamma_1 < \gamma_2 < \cdots$, contradicting $(**)$.

When these conditions are satisfied we say that $\Gamma$ has the \emph{ascending chain
condition}\index{ascending chain condition} (a.c.c.), or the \emph{maximal condition}. Reversing the order we can define
the \emph{descending chain condition}\index{descending chain condition} (d.c.c.), or \emph{minimal condition} in the same way.

If the set of ideals of a ring $A$ has the a.c.c., we say that $A$ is a \emph{Noetherian
ring}\index{Noetherian ring}, and if it has the d.c.c., that $A$ is an \emph{Artinian ring}.\index{Artinian ring} If $A$ is Noetherian (or
Artinian) and $B$ is a quotient of $A$ then $B$ has the same property; this is
obvious, since the set of ideals of $B$ is order-isomorphic to a subset of
that of $A$.\footnote{The a.c.c. and d.c.c. were first used in a paper of Emmy Noether (1882--1935), \textit{Idealtheorie in Ringbereichen}, Math.\ Ann., 83 (1921). Emil Artin (1898--1962) was, together with Emmy Noether, one of the founders of modern abstract algebra. As well as studying non-commutative rings whose one-sided ideals satisfy the d.c.c., he also discovered the Artin--Rees lemma, which will turn up in \S\ref{sec:8}.}

In the same way, we say that a module is \emph{Noetherian} or \emph{Artinian} if its set of
submodules satisfies the a.c.c. or the d.c.c. If $M$ has either of these properties,
then so do both its quotient modules and its submodules. (A subring of a
Noetherian or Artinian ring does not necessarily have the same property:
why not?)

A ring $A$ is Noetherian if and only if every ideal of $A$ is finitely generated.
(Proof, `only if': given an ideal $I$, consider a maximal element of the set of
finitely generated ideals contained in $I$; this must coincide with $I$. `If': given
an ascending chain $I_1 \subset I_2 \subset \cdots$ of ideals, $\bigcup I_n$ is also an ideal, so that by
assumption it can be generated by finitely many elements $a_1, \ldots, a_l$. There is
some $I_n$ which contains all the $a_i$, and the chain must stop there.)

In exactly the same way, an $A$-module $M$ is Noetherian if and only if
every submodule of $M$ is finitely generated. In particular $M$ itself must be
finitely generated, and if $A$ is Noetherian then this is also sufficient. Thus we
have the well-known fact that finite modules over a Noetherian ring are
Noetherian; we now give a proof of this in a more general form.

\begin{theorem}\label{thm:3.1}
Let $A$ be a ring and $M$ an $A$-module.
\begin{enumerate}[(i)]
  \item Let $M' \subset M$ be a submodule and $\varphi\colon M \to M/M'$ the natural map. If
        $N_1$ and $N_2$ are submodules of $M$ such that $N_1 \subset N_2$, $N_1 \cap M' = N_2 \cap M'$
        and $\varphi(N_1) = \varphi(N_2)$ then $N_1 = N_2$.
  \item Let $0 \to M' \to M \to M'' \to 0$ be an exact sequence of $A$-modules;
        if $M'$ and $M''$ are both Noetherian (or both Artinian), then so is $M$.
  \item Let $M$ be a finite $A$-module; then if $A$ is Noetherian (or Artinian),
        so is $M$.
\end{enumerate}
\end{theorem}

\begin{proof}
(i) is easy, and we leave it to the reader.

(ii) is obtained by applying (i) to an ascending (respectively descending)
chain of submodules of $M$.

(iii) If $M$ is generated by $n$ elements then it is a quotient of the free module
$A^n$, so that it is enough to show that $A^n$ is Noetherian (respectively Artinian).
However, this is clear from (ii) by induction on $n$.
\end{proof}

For a module $M$, it is equivalent to say that $M$ has both the a.c.c. and
the d.c.c., or that $M$ has finite length. Indeed, if $l(M) < \infty$ then $l(M_1) < l(M_2)$
for any two distinct submodules $M_1 \subset M_2 \subset M$, so that the two chain
conditions are clear. Conversely, if $M$ has the d.c.c. then we let $M_1$ be a
minimal non-zero submodule of $M$, let $M_2$ be a minimal element among all
submodules of $M$ strictly containing $M_1$, and proceed in the same way to
obtain an ascending chain $0 = M_0 \subset M_1 \subset M_2 \subset \cdots$; if $M$ also has the a.c.c.
then this chain must stop by arriving at $M$, so that $M$ has a composition
series.

Every submodule of the $\Z$-module $\Z$ is of the form $n\Z$, so that $\Z$ is
Noetherian, but not Artinian. Let $p$ be a prime, and write $W$ for the $\Z$-
module of rational numbers whose denominator is a power of $p$; then
the $\Z$-module $W/\Z$ is not Noetherian, but it is Artinian, since every proper
submodule of $W/\Z$ is either $0$ or is generated by $p^{-n}$ for $n = 1, 2, \ldots$ . This
shows that the a.c.c. and d.c.c. for modules are independent conditions,
but this is not the case for rings, as shown by the following result.

\begin{theorem}[Y.\ Akizuki]\label{thm:3.2}
An Artinian ring is Noetherian.
\end{theorem}

\begin{proof}
Let $A$ be an Artinian ring. It is sufficient to prove that $A$ has finite
length as an $A$-module. First of all, $A$ has only finitely many maximal ideals.
Indeed, if $p_1, p_2, \ldots$ is an infinite set of distinct maximal ideals then it is easy
to see that $p_1 \supset p_1p_2 \supset p_1p_2p_3 \supset \cdots$ is an infinite descending chain of ideals,
which contradicts the assumption. Thus, we let $p_1, p_2, \ldots, p_r$ be all the
maximal ideals of $A$ and set $I = p_1p_2 \cdots p_r = \rad(A)$. The descending chain
$I \supset I^2 \supset \cdots$ stops after finitely many steps, so that there is an $s$ such that
$I^s = I^{s+1}$. If we set $(0:I^s) = J$ then
\[
  (J:I) = ((0:I^s):I) = (0:I^{s+1}) = J;
\]
let's prove that $J = A$. By contradiction, suppose that $J \neq A$; then there exists
an ideal $J'$ which is minimal among all ideals strictly bigger than $J$. For any
$x \in J' - J$ we have $J' = Ax + J$. Now $I = \rad(A)$ and $J \neq J'$, so that by NAK
$J' \neq Ix + J$, and hence by minimality of $J'$ we have $Ix + J = J$, and this gives
$Ix \subset J$. Thus $x \in (J:I) = J$, which is a contradiction. Therefore $J = A$, so that
$I^s = 0$. Now consider the chain of ideals
\[
  A \supset p_1 \supset p_1p_2 \supset \cdots \supset p_1 \cdots p_r \supset (p_1 \cdots p_r)p_1 \supset \cdots \supset (p_1 \cdots p_r)^s = 0.
\]
Let $M$ and $Mp_i$ be any two consecutive terms in this chain; then $M/Mp_i$ is a
vector space over the field $A/p_i$, and since it is Artinian, it must be finite-
dimensional. Hence, $l(M/Mp_i) < \infty$, and therefore the sum $l(A)$ of these
terms is also finite.
\end{proof}

\begin{remark}
This theorem is sometimes referred to as Hopkins' theorem, but it was proved in the
above form by Akizuki [2] in 1935. It was rediscovered four years later by Hopkins [1], and
he proved it for non-commutative rings (a left-Artinian ring with unit is also left-Noetherian).
\end{remark}

\begin{theorem}\label{thm:3.3}
If $A$ is Noetherian then so are $A[X]$ and $A[[X]]$.
\end{theorem}

\begin{proof}
The statement for $A[X]$ is the well-known Hilbert basis theorem
(see, for example Lang, \textit{Algebra}, or [AM], p.~81), and we omit the proof.
We now briefly run through the proof for $A[[X]]$. Set $B = A[[X]]$, and
let $I$ be an ideal of $B$; we will prove that $I$ is finitely generated. Write
$I(r)$ for the ideal of $A$ formed by the leading coefficients $a_r$ of
$f = a_rX^r + a_{r+1}X^{r+1} + \cdots$ as $f$ runs through $I \cap X^rB$; then we have
\[
  I(0) \subset I(1) \subset I(2) \subset \cdots .
\]
Since $A$ is Noetherian, there is an $s$ such that $I(s) = I(s+1) = \cdots$; moreover,
each $I(i)$ is finitely generated. For each $i$ with $0 \le i \le s$ we take finitely many
elements $a_{i\nu} \in A$ generating $I(i)$, and choose $g_{i\nu} \in I \cap X^iB$ having $a_{i\nu}$ as the
coefficient of $X^i$. These $g_{i\nu}$ now generate $I$. Indeed, for $f \in I$ we can take
a linear combination $g_0$ of the $g_{0\nu}$ with coefficients in $A$ such that
$f - g_0 \in I \cap XB$, then take a linear combination $g_1$ of the $g_{1\nu}$ with
coefficients in $A$ such that $f - g_0 - g_1 \in I \cap X^2B$, and proceeding in
the same way we get
\[
  f - g_0 - g_1 - \cdots - g_s \in I \cap X^{s+1}B.
\]
Now $I(s+1) = I(s)$, so we can take a linear combination $g_{s+1}$ of the $Xg_{s\nu}$
with coefficients in $A$ such that
\[
  f - g_0 - g_1 - \cdots - g_{s+1} \in I \cap X^{s+2}B.
\]
We now proceed in the same way to get $g_{s+2}, \ldots$. For $i \le s$, each $g_i$
is a linear combination of the $g_{i\nu}$ with coefficients in $A$, and, for
$i > s$, a combination of the elements $X^{i-s}g_{s\nu}$. For each $i > s$ we write
$g_i = \sum_\nu a_{i\nu}X^{i-s}g_{s\nu}$, and then for each $\nu$ we set
$h_\nu = \sum_{i \ge s} a_{i\nu}X^{i-s}$; $h_\nu$ is an element of $B$, and
\[
  f = g_0 + \cdots + g_{s-1} + \sum_\nu h_\nu g_{s\nu}.
\]
\end{proof}

A ring $A[b_1, \ldots, b_n]$ which is finitely generated as a ring over a
Noetherian ring $A$ is a quotient of a polynomial ring $A[X_1, \ldots, X_n]$, and
so by the Hilbert basis theorem is again Noetherian. We now give some
other criteria for a ring to be Noetherian.

\begin{theorem}[I.\,S.\ Cohen]\label{thm:3.4}
If all the prime ideals of a ring $A$ are finitely
generated then $A$ is Noetherian.
\end{theorem}

\begin{proof}
Write $\Gamma$ for the set of ideals of $A$ which are not finitely generated.
If $\Gamma \neq \varnothing$ then by Zorn's lemma $\Gamma$ contains a maximal element $I$. Then $I$ is
not a prime ideal, so that there are elements $x, y \in A$ with $x \notin I$, $y \notin I$ but
$xy \in I$. Now $I + Ay$ is bigger than $I$, and hence is finitely generated, so
that we can choose $u_1, \ldots, u_n \in I$ such that
\[
  I + Ay = (u_1, \ldots, u_n, y).
\]
Moreover, $I:y = \{a \in A \mid ay \in I\}$ contains $x$, and is thus bigger than
$I$, so it has a finite system of generators $(v_1, \ldots, v_m)$. Finally, it is
easy to check that $I = (u_1, \ldots, u_n, yu_1, \ldots, yu_n, v_1y, \ldots, v_my)$; hence, $I \notin \Gamma$, which is a
contradiction. Therefore $\Gamma = \varnothing$.
\end{proof}

\begin{theorem}\label{thm:3.5}
Let $A$ be a ring and $M$ an $A$-module. Then if $M$ is a
Noetherian module, $A/\ann(M)$ is a Noetherian ring.
\end{theorem}

\begin{proof}
If we set $\bar{A} = A/\ann(M)$ and view $M$ as an $\bar{A}$-module, then the
submodules of $M$ as an $A$-module or $\bar{A}$-module coincide, so that $M$ is
also Noetherian as an $\bar{A}$-module. We can thus replace $A$ by $\bar{A}$, and then
$\ann(M) = (0)$. Now letting $M = A\omega_1 + \cdots + A\omega_n$, we can embed $A$ in $M^n$
by means of the map $a \mapsto (a\omega_1, \ldots, a\omega_n)$. By Theorem 1, $M^n$ is a Noetherian
module, so that its submodule $A$ is also Noetherian. (This theorem can
be expressed by saying that a ring having a faithful Noetherian module
is Noetherian.)
\end{proof}

\begin{theorem}[E.\ Formanek~\cite{Fo1}]\label{thm:3.6}
Let $A$ be a ring, and $B$ an $A$-module
which is finitely generated and faithful over $A$. Assume that the set of
submodules of $B$ of the form $IB$ with $I$ an ideal of $A$ satisfies the a.c.c.;
then $A$ is Noetherian.
\end{theorem}

\begin{proof}
It will be enough to show that $B$ is a Noetherian $A$-module. By
contradiction, suppose that it is not; then the set
\[
  \{IB \mid I \text{ is an ideal of } A \text{ and } B/IB \text{ is non-Noetherian as } A\text{-module}\}
\]
contains $(0)$ and so is non-empty, so that by assumption it contains a
maximal element. Let $IB$ be one such maximal element; then replacing $B$
by $B/IB$ and $A$ by $A/\ann(B/IB)$ we see that we can assume that $B$ is a
non-Noetherian $A$-module, but for any non-zero ideal $I$ of $A$ the quotient
$B/IB$ is Noetherian.

Next we set
\[
  \mathcal{I} = \{N \mid N \text{ is a submodule of } B \text{ and } B/N \text{ is a faithful } A\text{-module}\}.
\]
If $B = Ab_1 + \cdots + Ab_n$ then for a submodule $N$ of $B$,
\[
  N \in \mathcal{I} \iff \forall a \in A - \{0\}, \quad \{ab_1, \ldots, ab_n\} \nsubseteq N.
\]
From this, one sees at once that Zorn's lemma applies to $\mathcal{I}$; hence there
exists a maximal element $N_0$ of $\mathcal{I}$. If $B/N_0$ is Noetherian then $A$ is a
Noetherian ring, and thus $B$ is Noetherian, which contradicts our
hypothesis. It follows that on replacing $B$ by $B/N_0$ we arrive at a module
$B$ with the following properties:
\begin{enumerate}[(1)]
  \item $B$ is non-Noetherian as an $A$-module;
  \item for any ideal $I \neq (0)$ of $A$, $B/IB$ is Noetherian;
  \item for any submodule $N \neq (0)$ of $B$, $B/N$ is not faithful as an $A$-module.
\end{enumerate}
Now let $N$ be any non-zero submodule of $B$. By (3) there is an element
$a \in A$ with $a \neq 0$ such that $a(B/N) = 0$, that is such that $aB \subset N$. By (2)
$B/aB$ is a Noetherian module, so that $N/aB$ is finitely generated; but since
$B$ is finitely generated so is $aB$, and hence $N$ itself is finitely generated.
Thus, $B$ is a Noetherian module, which contradicts (1).
\end{proof}

As a corollary of this theorem we get the following result.

\begin{theorem}\label{thm:3.7}
\begin{enumerate}[(i)]
  \item (Eakin--Nagata theorem). Let $B$ be a Noetherian ring, and $A$ a subring
        of $B$ such that $B$ is finite over $A$; then $A$ is also a Noetherian ring.
  \item Let $B$ be a non-commutative ring whose right ideals have the a.c.c.,
        and let $A$ be a commutative subring of $B$. If $B$ is finitely generated as a
        left $A$-module then $A$ is a Noetherian ring.
  \item Let $B$ be a non-commutative ring whose two-sided ideals have the
        a.c.c., and let $A$ be a subring contained in the centre of $B$; if $B$ is finitely
        generated as an $A$-module then $A$ is a Noetherian ring.
\end{enumerate}
\end{theorem}

\begin{proof}
$B$ has a unit, so is faithful as an $A$-module. Hence it is enough to
apply the previous theorem.
\end{proof}

\begin{remark}
Part (i) of Theorem 7 was proved in Eakin's thesis [1] in 1968,
and the same result was obtained independently by Nagata [9] a little
later. Subsequently many alternative proofs and extensions to the non-
commutative case were published; the most transparent of these seems to
be Formanek's result~\cite{Fo1}, which we have given above in the form of
Theorem 6. However, this also goes back to the idea of the proofs of
Eakin and Nagata.
\end{remark}

% ============ 习题 ============
\section*{Exercises to \S\ref{sec:3}}
\emph{Prove the following propositions.}

\begin{exercise}\label{ex:3.1-ex}
Let $I_1, \ldots, I_n$ be ideals of a ring $A$ such that $I_1 \cap \cdots \cap I_n = (0)$; if each $A/I_i$ is a
Noetherian ring then so is $A$.
\end{exercise}

\begin{exercise}\label{ex:3.2-ex}
Let $A$ and $B$ be Noetherian rings, and $f\colon A \to C$ and $g\colon B \to C$ ring
homomorphisms. If both $f$ and $g$ are surjective then the fibre product
$A \times_C B$ (that is, the subring of the direct product $A \times B$ given by
$\{(a, b) \in A \times B \mid f(a) = g(b)\}$) is a Noetherian ring.
\end{exercise}

\begin{exercise}\label{ex:3.3-ex}
Let $A$ be a local ring such that the maximal ideal $m$ is principal and
$\bigcap_{n=1}^{\infty} m^n = (0)$. Then $A$ is Noetherian, and every non-zero ideal of $A$ is a
power of $m$.
\end{exercise}

\begin{exercise}\label{ex:3.4-ex}
Let $A$ be an integral domain with field of fractions $K$. A \emph{fractional ideal} $I$ of
$A$ is an $A$-submodule $I$ of $K$ such that $I \neq 0$ and $aI \subset A$ for some $0 \neq a \in K$.
The product of two fractional ideals is defined in the same way as the
product of two ideals. If $I$ is a fractional ideal of $A$ we set $I^{-1} = \{\alpha \in K \mid \alpha I
\subset A\}$; this is also a fractional ideal, and $II^{-1} \subset A$. In the particular case
that $II^{-1} = A$ we say that $I$ is \emph{invertible}. An invertible fractional ideal of $A$
is finitely generated as an $A$-module.
\end{exercise}

\begin{exercise}\label{ex:3.5-ex}
If $A$ is a UFD, the only ideals of $A$ which are invertible as fractional ideals
are the principal ideals.
\end{exercise}

\begin{exercise}\label{ex:3.6-ex}
Let $A$ be a Noetherian ring, and $\varphi\colon A \to A$ a homomorphism of rings.
Then if $\varphi$ is surjective it is also injective, and hence an automorphism of $A$.
\end{exercise}

\begin{exercise}\label{ex:3.7-ex}
If $A$ is a Noetherian ring then any finite $A$-module is of finite presentation,
but if $A$ is non-Noetherian then $A$ must have finite $A$-modules which are
not of finite presentation.
\end{exercise}
